\section{Private Data Dissemination and Retrieval}
\label{s:gossip}
Since private data is excluded from the consensus-driven block delivery, we utilize a specialized peer-to-peer protocol to ensure data availability and fault tolerance. The protocol operates in two complementary phases: proactive dissemination during simulation, and reactive retrieval during commit.

\subsection{Proactive Dissemination}
\label{s:dissemination}
During the simulation phase, each endorsing node stores the generated private write-sets in its local transient store and proactively pushes them to a subset of authorized remote peers via a peer-to-peer protocol. For each collection involved in the transaction, the endorser sends the entire collection's private write-set as a single message to eligible peers. The set of target peers is determined by evaluating the collection's membership policy $\mathcal{M}$ against the current view of active peers in the network. The number of peers targeted is governed by $n_{dissem}$, and the endorser waits until at least $n_{required}$ peers acknowledge receipt before returning the signed endorsement response to the client. This eager strategy ensures that the private data is already present on most authorized nodes by the time the corresponding block is received, minimizing the need for reactive pulling during the commit phase. If the endorser cannot obtain $n_{required}$ acknowledgments before the timeout, it returns a dissemination failure to the client; the client may then retry with a different set of endorsers or abort the transaction. The endorser itself does not re-attempt dissemination, as the client is responsible for orchestrating retries---this keeps the endorser stateless with respect to dissemination attempts. Algorithm~\ref{alg:dissemination} summarizes the protocol.

\begin{algorithm}[h]
\centering
\small
\begin{tabular}{|p{0.92\columnwidth}|}
\hline
\textbf{Proactive Dissemination Protocol} \\
\hline
\textbf{Input:} Transaction $tx$ with private write-sets \\[3pt]
\textbf{At Endorsing Node} $e$ (after simulation): \\
\quad 1. Store all private write-sets for $tx$ in local transient store \\
\quad 2. \textbf{for each} collection $c$ in $tx$'s private write-sets \textbf{do} \\
\quad\quad\quad $W_c \leftarrow$ collection $c$'s private write-set \\
\quad\quad\quad $P \leftarrow \text{SelectAuthorizedPeers}(c.\mathcal{M}, c.n_{dissem})$ \\
\quad\quad\quad \textbf{for each} peer $p_i \in P$ \textbf{do} \\
\quad\quad\quad\quad\quad Push $(tx.id, c.id, W_c)$ to $p_i$ via gossip \\
\quad\quad\quad Wait until $|\{p_i : \text{ACK received}\}| \geq c.n_{required}$ \\
\quad\quad\quad\quad\quad or timeout expires \\
\quad\quad\quad \textbf{if} timeout \textbf{then} return dissemination failure \\
\quad 3. Return signed endorsement response to client \\[3pt]
\textbf{At Receiving Peer} $p_i$: \\
\quad 1. Verify sender $e$ is authorized for collection $c$ \\
\quad 2. Verify ledger height of $tx$ is within acceptable range \\
\quad 3. Store $(tx.id, c.id, W_c)$ in local transient store \\
\quad 4. Send ACK to endorser $e$ \\
\hline
\end{tabular}
\caption{Proactive dissemination protocol executed during the simulation phase. The endorser pushes each collection's full private write-set to authorized peers before returning the endorsement.}
\label{alg:dissemination}
\end{algorithm}


\subsection{Reactive Retrieval}
\label{s:retrieval}
Despite hash pre-images can be disseminated to target peers before they receive the block that contains the corresponding hashes, a robust solution is needed with withstand network partitions and peer crashes which prevent the hash pre-images from reaching all the required peers. Furthermore, if new peers are added to the Blockchain, they can only obtain the required pre-images by pulling them from peers that posses them.

To that end, prior to committing a block to the ledger, a peer iterates through all hashed key writes in all transactions of the Block and analyzes whether it needs to pull private data for each hashes key write. A hashed key write is considered as a viable candidate for pulling pre-images if:
\begin{itemize}
	\item The transaction passed validation of endorsement policies {Are we into abstracting out EPs? If yes then how should I phrase this?}. This is needed in order to prevent a denial-of-service attack in which a malicious party creates  transactions with hashed writes for which there are no pre-image in any peer in the network. Peers would try to pull these transactions and fail, hence considerably slowing down throughput.
	\item Evaluation of the collection policy on the peer's identity returns "true" {Once we define collection policies more abstractly we need to change this as well}: this is required since peers that receive pre-image requests for a certain hashed write also evaluate the same collection policy on the requester peer's identity in order to determine if that peer is part of the collection the hashed write is associated with.
	\item Whether the peer is missing the pre-image from its transient store, since by definition this means the pre-image hasn't been disseminated to the peer before the commit of the block.
	\item {Should we mention purged private data?}
\end{itemize}

Once the peer decides it needs to pull a hashed write, it needs to select potential peers which might have the pre-image it needs, and send them requests for the digests.
Upon reception of responses containing pre-images, the peer validates that the hashes on the pre-images received matches the hashed write in the block. It follows from the collision resistance property of the cryptographic hash function that if the responding peer sent a pre-image of a hashed write, then with overwhelming probability, the pre-image is identical to all pre-images in all peers in the network, hence consistency is preserved.

When sending requests for pre-images, a requester peer doesn't know up-front whether the responder peer has the pre-image, and every such "wrong guess" results in a retry which delays the block commit and thus hinders performance.

For this reason, the peers prioritize requesting a pre-image from the peers which endorsed the transaction {This is very Fabric specific, so we might want to remove this?} from which the hashed write originated. This is possible because peers constantly maintain a membership view of alive peers by periodically disseminating heartbeats. 



\subsection{Dynamic Membership and Reconciliation}
\label{s:dynamic-membership}
Collection membership may change over the lifetime of a network. When a new member is added to a collection's membership policy $\mathcal{M}$, the peer detects the configuration change during the commit of the block containing the policy update and initiates a retroactive reconciliation process. Specifically, the peer scans the committed blocks for historical transactions that wrote to the collection and marks them as missing eligible data. A background reconciliation service then periodically pulls these historical private write-sets from existing collection members using the same hash-based verification as the reactive retrieval protocol.

When a member is removed from a collection, the removed peer retains access only to data committed before its removal. The updated membership policy excludes it from both proactive dissemination and reactive retrieval eligibility checks, ensuring that future private data is not shared with the departing organization.

\subsection{Consistency and Fault Tolerance}
A critical property of the retrieval mechanism is its idempotency: since multiple peers may attempt to pull the same missing data concurrently, the system utilizes the versioned hashes on the public ledger as unique identifiers, ensuring that nodes eventually converge on a consistent Private State even in the presence of network partitions or concurrent block commits.

\smallskip
\noindent\textbf{Remark} (Integrity Independence). \textit{Data integrity is preserved even when peers attempt to provide fraudulent pre-images, because the requester validates every incoming data piece against the on-chain hash recorded during the consensus phase. This integrity guarantee operates independently of the CFT assumption described in \S\ref{pdc}: while the \textbf{liveness} of dissemination depends on honest peers following the protocol, the \textbf{correctness} of retrieved data is ensured solely by cryptographic hash verification. A single honest data source is sufficient for correctness.}

